\documentclass[handout]{beamer}
\mode<presentation>
\usetheme[secheader]{Boadilla}
\usecolortheme{whale}

\usepackage[utf8]{inputenc}
\usepackage[T1]{fontenc}
\usepackage[indonesian]{babel}
\usepackage{lmodern}
\usepackage{graphicx}
\usepackage{bookmark}

\setbeamertemplate{footline}
{
	\leavevmode%
	\hbox{%
		\begin{beamercolorbox}[wd=.333333\paperwidth,ht=2.25ex,dp=1ex,center]{author in head/foot}%
			\usebeamerfont{author in head/foot}\insertshortauthor%
		\end{beamercolorbox}%
		\begin{beamercolorbox}[wd=.433333\paperwidth,ht=2.25ex,dp=1ex,center]{title in head/foot}%
			\usebeamerfont{title in head/foot}\insertshorttitle
		\end{beamercolorbox}%
		\begin{beamercolorbox}[wd=.233333\paperwidth,ht=2.25ex,dp=1ex,right]{date in head/foot}%
			\usebeamerfont{date in head/foot}\insertshortdate{}\hspace*{2em} \insertframenumber{} / \inserttotalframenumber\hspace*{2ex}
	\end{beamercolorbox}}%
	\vskip0pt%
}

\title[Bab 15 -- Integrasi \& Deployment]{Pemrograman Web}
\subtitle{Bab 15: Integrasi API, Pengujian, Keamanan, Deployment Dasar, Administrasi Web}
\author{Cecep Suwanda, S.Si., M.Kom.}
\institute{Universitas Bale Bandung}
\date{\today}

\begin{document}

% Slide Judul
\begin{frame}
	\titlepage
\end{frame}

% Outline dan CPMK
\begin{frame}{Outline Bab 15}
	\begin{block}{Sub-CPMK}
		\begin{itemize}
			\item \textbf{3.3}: Mengintegrasikan API, front-end $\leftrightarrow$ back-end
			\item \textbf{4.1}: Pengujian dan evaluasi performa
			\item \textbf{4.2}: Keamanan web (validasi, prepared statement, XSS, CSRF)
			\item \textbf{4.3}: Deployment ke produksi dan administrasi
		\end{itemize}
	\end{block}
	\begin{itemize}
		\item Integrasi API dan Fetch
		\item Pengujian (fungsional, kompatibilitas, performa)
		\item Keamanan (validasi, SQL injection, XSS, CSRF, session)
		\item Deployment dan Administrasi (dev vs prod, php.ini, backup, CI/CD)
		\item Referensi Standar
	\end{itemize}
\end{frame}

% ========== INTEGRASI API ==========
\begin{frame}{Integrasi API dan Pertukaran Data}
	\begin{itemize}
		\item \textbf{API}: antarmuka untuk pertukaran data; arsitektur \textbf{REST}; format \textbf{JSON}.
		\item \textbf{Fetch API}: HTTP request dari JavaScript; berbasis Promise; async/await; tidak reload halaman (AJAX).
		\item Alur: Front-End (JS) $\to$ HTTP Request $\to$ API Endpoint (PHP) $\to$ SQL $\to$ MySQL $\to$ JSON Response $\to$ Front-End.
		\item Memisahkan front-end dan back-end; mendukung multi-platform (web, mobile).
	\end{itemize}
\end{frame}

% ========== PENGUJIAN ==========
\begin{frame}{Pengujian Aplikasi Web}
	\begin{itemize}
		\item Memverifikasi aplikasi berfungsi sesuai harapan, bebas bug, memenuhi standar.
		\item \textbf{Shift-left}: pengujian diintegrasikan di setiap tahap, bukan hanya di akhir.
	\end{itemize}
	\begin{block}{Tujuan}
		Verifikasi fungsionalitas, deteksi bug, keamanan, performa, kompatibilitas browser.
	\end{block}
\end{frame}

\begin{frame}{Jenis Pengujian}
	\begin{tabular}{|l|l|l|}
		\hline
		\textbf{Jenis} & \textbf{Fokus} & \textbf{Alat} \\
		\hline
		Fungsional & Fitur sesuai spesifikasi & Selenium, Cypress \\
		\hline
		Kompatibilitas & Browser, OS, device & BrowserStack \\
		\hline
		Performa & Kecepatan, responsivitas & Lighthouse, k6 \\
		\hline
		Keamanan & Celah kerentanan & OWASP ZAP \\
		\hline
	\end{tabular}
\end{frame}

\begin{frame}{DevTools: Elements, Console, Network, Lighthouse}
	\begin{itemize}
		\item \textbf{Elements}: inspeksi HTML/CSS, cek layout
		\item \textbf{Console}: log, error JS, eksekusi kode
		\item \textbf{Network}: HTTP request/response, status, waktu; penting untuk debug API
		\item \textbf{Lighthouse}: audit performa, aksesibilitas, SEO, best practices
	\end{itemize}
\end{frame}

\begin{frame}{Strategi Pengujian}
	\begin{block}{Manual vs Otomatis}
		Manual: eksplorasi, usabilitas. Otomatis: regresi, CI/CD.
	\end{block}
	\begin{block}{Piramida: Unit $\to$ Integrasi $\to$ E2E}
		Unit (banyak, cepat) $\to$ Integrasi (sedang) $\to$ E2E (sedikit, lambat).
	\end{block}
\end{frame}

\begin{frame}{Pengujian Fungsional}
	\begin{itemize}
		\item \textbf{Skenario positif}: input valid $\to$ output benar
		\item \textbf{Skenario negatif}: input tidak valid $\to$ pesan error jelas
		\item \textbf{Boundary}: nilai di batas min/max
		\item Validasi form di server (PHP) wajib; validasi client (JS) hanya untuk UX.
	\end{itemize}
\end{frame}

\begin{frame}{Metrik Performa: Core Web Vitals}
	\begin{tabular}{|l|l|l|}
		\hline
		\textbf{Metrik} & \textbf{Target baik} & \textbf{Deskripsi} \\
		\hline
		FCP & $\leq$ 1,8 s & Konten pertama muncul \\
		\hline
		LCP & $\leq$ 2,5 s & Elemen terbesar dirender \\
		\hline
		INP & $\leq$ 200 ms & Respons interaksi \\
		\hline
		CLS & $\leq$ 0,1 & Stabilitas layout \\
		\hline
	\end{tabular}
	\vspace{0.5em}
	Alat: \textbf{Lighthouse} (DevTools), PageSpeed Insights, WebPageTest.
\end{frame}

% ========== KEAMANAN ==========
\begin{frame}{Keamanan Aplikasi Web}
	\begin{block}{Prinsip}
		\textit{Jangan pernah percaya input dari client}. Semua data dari browser harus divalidasi.
	\end{block}
	\begin{itemize}
		\item \textbf{CIA Triad}: Confidentiality, Integrity, Availability
		\item Ancaman umum: SQL Injection, XSS, CSRF, Session Hijacking, Brute Force
		\item \textbf{Defence in depth}: keamanan berlapis, bukan satu titik
	\end{itemize}
\end{frame}

\begin{frame}{Validasi vs Sanitasi}
	\begin{block}{Validasi}
		Memeriksa format, panjang, rentang. Tolak data tidak valid. Wajib di server.
	\end{block}
	\begin{block}{Sanitasi}
		Membersihkan karakter berbahaya (\texttt{htmlspecialchars}, \texttt{trim}). Sebelum simpan/tampilkan.
	\end{block}
	Validasi client (JS) hanya untuk UX; validasi server (PHP) adalah pertahanan utama.
\end{frame}

\begin{frame}{Prepared Statement: Mencegah SQL Injection}
	\begin{itemize}
		\item Query template dengan placeholder \texttt{?}; data dikirim terpisah.
		\item Database memperlakukan data sebagai literal, bukan perintah SQL.
		\item Contoh: \texttt{prepare("SELECT * FROM users WHERE email = ?")}, \texttt{bind\_param("s", \$email)}.
		\item Jangan gabungkan input langsung ke query string.
	\end{itemize}
\end{frame}

\begin{frame}{Session: Pengelolaan Aman}
	\begin{itemize}
		\item \texttt{session\_start()}; simpan data di \texttt{\$\_SESSION}.
		\item \textbf{Regenerasi ID} setelah login: \texttt{session\_regenerate\_id(true)} (cegah session fixation).
		\item Logout: hapus \texttt{\$\_SESSION}, hancurkan cookie.
		\item Konfigurasi: \texttt{cookie\_httponly}, \texttt{cookie\_secure}, \texttt{cookie\_samesite}, \texttt{use\_strict\_mode}.
	\end{itemize}
\end{frame}

\begin{frame}{Cookie: HttpOnly, Secure, SameSite}
	\begin{tabular}{|l|l|l|}
		\hline
		\textbf{Atribut} & \textbf{Fungsi} & \textbf{Melindungi} \\
		\hline
		HttpOnly & Cookie tidak bisa diakses JS & XSS (pencurian cookie) \\
		\hline
		Secure & Hanya via HTTPS & Penyadapan \\
		\hline
		SameSite=Strict & Tidak dikirim lintas situs & CSRF \\
		\hline
	\end{tabular}
\end{frame}

\begin{frame}{Pencegahan XSS: Output Encoding}
	\begin{itemize}
		\item XSS: penyerang sisipkan skrip JS yang dieksekusi di browser korban.
		\item Tipe: Stored (di DB), Reflected (dari URL), DOM-based.
		\item \textbf{Pencegahan}: \texttt{htmlspecialchars(\$data, ENT\_QUOTES, 'UTF-8')} setiap output dari input pengguna.
		\item Prinsip: sanitasi saat output, bukan saat input.
	\end{itemize}
\end{frame}

\begin{frame}{Pencegahan CSRF: Token}
	\begin{itemize}
		\item CSRF: pengguna login tanpa sadar mengirim request berbahaya dari situs penyerang.
		\item \textbf{CSRF token}: nilai unik di form; divalidasi di server.
		\item Generate: \texttt{bin2hex(random\_bytes(32))}; validasi: \texttt{hash\_equals()}.
		\item Cookie \texttt{SameSite=Strict} juga membantu cegah CSRF.
	\end{itemize}
\end{frame}

\begin{frame}{Load Testing dan Optimasi Performa}
	\begin{itemize}
		\item \textbf{Load testing}: mengukur perilaku di bawah beban; alat: JMeter, Locust, k6.
		\item \textbf{Optimasi}: caching (browser, server, opcode), minifikasi CSS/JS, lazy loading gambar.
		\item \texttt{loading="lazy"} pada \texttt{<img>}; gzip/brotli; Cache-Control; WebP.
	\end{itemize}
\end{frame}

\begin{frame}{OWASP Top 10}
	\small
	\begin{tabular}{|l|l|}
		\hline
		\textbf{Ancaman} & \textbf{Mitigasi} \\
		\hline
		Injection (SQL, dll.) & Prepared statement \\
		\hline
		Broken Access Control & Cek izin di server \\
		\hline
		Cryptographic Failures & HTTPS, \texttt{password\_hash} \\
		\hline
		Security Misconfig & Hardening, non-default password \\
		\hline
		XSS & Output encoding \\
		\hline
		... & Lakukan audit berkala \\
		\hline
	\end{tabular}
\end{frame}

\begin{frame}{Session Fixation dan Hijacking}
	\begin{block}{Session Fixation}
		Penyerang paksa korban pakai session ID penyerang. \textbf{Mitigasi}: \texttt{session\_regenerate\_id(true)} setelah login.
	\end{block}
	\begin{block}{Session Hijacking}
		Pencurian session ID (via XSS, penyadapan). \textbf{Mitigasi}: HttpOnly, Secure, SameSite, timeout session.
	\end{block}
\end{frame}

% ========== DEPLOYMENT ==========
\begin{frame}{Development vs Production}
	\begin{tabular}{|l|l|l|}
		\hline
		\textbf{Aspek} & \textbf{Dev} & \textbf{Prod} \\
		\hline
		\texttt{display\_errors} & On & Off \\
		\hline
		\texttt{log\_errors} & Off & On \\
		\hline
		Database & Lokal, dummy & Server, data asli \\
		\hline
		HTTPS & Tidak wajib & Wajib \\
		\hline
		Password & Sederhana & Kuat dan unik \\
		\hline
		Backup & Tidak perlu & Rutin, otomatis \\
		\hline
	\end{tabular}
\end{frame}

\begin{frame}{Deployment dan Siklus Rilis}
	\begin{enumerate}
		\item Develop $\to$ Test $\to$ Stage $\to$ Deploy $\to$ Monitor $\to$ Feedback
	\end{enumerate}
	\begin{block}{Checklist pra-deploy}
		Tes lulus; database migrasi; env var; \texttt{display\_errors} off; HTTPS; backup; \texttt{.env}, \texttt{.git/}, \texttt{vendor/} tidak bisa diakses publik.
	\end{block}
\end{frame}

\begin{frame}{Server Web: Apache vs Nginx}
	\begin{itemize}
		\item \textbf{Apache}: mod\_php, .htaccess; mudah konfigurasi.
		\item \textbf{Nginx}: event-driven, PHP-FPM; performa tinggi; populer sebagai reverse proxy.
	\end{itemize}
\end{frame}

\begin{frame}{Environment Variable dan .env}
	\begin{itemize}
		\item Simpan konfigurasi (DB, API key, mode debug) di luar kode.
		\item File \texttt{.env}: \textbf{jangan commit ke Git}; tambah ke \texttt{.gitignore}.
		\item Library \texttt{vlucas/phpdotenv} membaca \texttt{.env} $\to$ \texttt{\$\_ENV}.
	\end{itemize}
\end{frame}

\begin{frame}{Konfigurasi PHP untuk Produksi}
	\begin{itemize}
		\item \texttt{display\_errors = Off}; \texttt{log\_errors = On}
		\item \texttt{opcache.enable = 1} (performa)
		\item \texttt{expose\_php = Off}
		\item \texttt{session.cookie\_httponly}, \texttt{session.cookie\_secure}, \texttt{session.cookie\_samesite}
	\end{itemize}
\end{frame}

\begin{frame}{Upload, Permissions, Deploy}
	\begin{itemize}
		\item Metode: Git (\texttt{git clone}/\texttt{pull}) atau SFTP/SCP.
		\item Permission: folder \texttt{755}, file \texttt{644}, upload/cache \texttt{775}.
		\item Langkah: SSH, install stack, clone, \texttt{composer install}, buat \texttt{.env}, import DB, \texttt{chown www-data}, virtual host, HTTPS (Let's Encrypt).
	\end{itemize}
\end{frame}

\begin{frame}{Reverse Proxy dan HTTPS}
	\begin{itemize}
		\item \textbf{Reverse proxy} (Nginx): SSL termination, load balancing, caching, keamanan.
		\item \textbf{HTTPS}: enkripsi TLS; Let's Encrypt (sertifikat gratis); \texttt{certbot}.
	\end{itemize}
\end{frame}

\begin{frame}{Hardening Server}
	\begin{itemize}
		\item Firewall (UFW): hanya port 22, 80, 443
		\item Update rutin (OS, software)
		\item User non-root untuk aplikasi
		\item SSH: nonaktifkan root login; key-based auth
		\item Port 3306 (MySQL) diblokir dari publik
	\end{itemize}
\end{frame}

\begin{frame}{Backup dan Restore Database}
	\begin{itemize}
		\item \texttt{mysqldump} untuk backup; jadwalkan dengan \texttt{cron} (harian).
		\item Simpan di lokasi terpisah; uji restore secara berkala.
		\item Contoh cron: \texttt{0 2 * * *} (setiap hari jam 02:00).
	\end{itemize}
\end{frame}

\begin{frame}{CI/CD: Otomasi Build dan Deploy}
	\begin{itemize}
		\item \textbf{CI}: integrasi kode, tes otomatis (setiap push).
		\item \textbf{CD}: deploy otomatis setelah tes lulus.
		\item Alat: GitHub Actions, GitLab CI, Jenkins.
		\item Pipeline: Push $\to$ Build $\to$ Test $\to$ Deploy Staging $\to$ Deploy Prod
	\end{itemize}
\end{frame}

\begin{frame}{Pemantauan di Produksi}
	\begin{itemize}
		\item \textbf{Log}: error log PHP, access log server, slow query MySQL
		\item \textbf{Error tracking}: Sentry, Bugsnag (exception real-time)
		\item \textbf{Uptime}: UptimeRobot, Pingdom
	\end{itemize}
\end{frame}

% ========== REFERENSI ==========
\begin{frame}{Referensi Standar}
	\begin{tabular}{|l|l|}
		\hline
		\textbf{Sumber} & \textbf{Topik} \\
		\hline
		MDN Web Docs & HTML, CSS, JS, Web API \\
		\hline
		OWASP & Keamanan web, checklist \\
		\hline
		web.dev & Performa, Core Web Vitals \\
		\hline
		PHP The Right Way & Best practices PHP \\
		\hline
		PHP Manual & Dokumentasi resmi \\
		\hline
	\end{tabular}
\end{frame}

% ========== RANGKUMAN ==========
\begin{frame}{Rangkuman Bab 15}
	\begin{itemize}
		\item \textbf{Sub-CPMK 3.3}: API REST, Fetch API, JSON; front-end $\leftrightarrow$ back-end.
		\item \textbf{Sub-CPMK 4.1}: Pengujian fungsional, performa (FCP, LCP, INP, CLS); Lighthouse; load testing.
		\item \textbf{Sub-CPMK 4.2}: Validasi di server; prepared statement; output encoding (XSS); CSRF token; session/cookie aman; OWASP Top 10.
		\item \textbf{Sub-CPMK 4.3}: Dev vs prod; .env; php.ini; reverse proxy; HTTPS; hardening; backup; CI/CD; monitoring.
	\end{itemize}
\end{frame}

\begin{frame}{}
	\centering
	\vfill
	\textbf{\Large Pertanyaan?}
	\vfill
\end{frame}

\end{document}
